\paragraph{Cross-circuit automatic repair recovers, but does not exceed, strong native SPMV.}
Starting from the sealed SPMV R1 route, the generic current-STA selector chooses 6008/pin~12. Native repair and FULL/STA reduce CPD from 7.045986460951781 to 7.038347682453150 ns, improving two setup components without regressions. Fresh STA then identifies a relaxed opportunity at 6395/pin~2. The automatic conflict scope also releases 6394/pin~2; native recovery lowers CPD to 7.034697713237392 ns. Both exact policies accept both steps and end at the same physical state. Compared with initial R1, 49 fixed setup/domain components improve, none regresses, and no hold violation increases. Final resources are two SB pairs, 967 CB sources and planar wirelength 148591, versus one, 967 and 148617 initially. Thus the delay gain requires an additional SB pair, despite lower planar wirelength.

This final CPD equals the unmodified A0 value, not a gain over A0. The sealed A0 route uses zero SB pairs, 967 CB sources and planar wirelength 148463. Running the same worst-current-STA operator from A0 yields no improving capacity or relaxed-distance opportunity and neither policy accepts a change. R1 recovery therefore does not establish superiority to the stronger native starting state, and its final resource use remains higher. The R1 two-policy controller takes 376.123 s for two distinct transactions and four native connection-search attempts, with 0.296 s inside native transactions. Candidate audits take 56.718 s, candidate FULL/STA wrappers 72.044 s, and all other controller/preparation/query work is included in the total. A0's no-candidate diagnostic costs 82.888 s. These are measured incremental prototype costs, not complete initial-routing times or independent per-policy speedups.

The observed SPMV R1 and Tf states have identical physical routes, actual input identities, normalized native commands and complete endpoint vectors. The Tf controller therefore explicitly reuses the completed R1 transactions, FULL/STA and matching query work, reaching the same states with zero new native repair core work. Its 51.354 s incremental reuse/selection overhead is reported separately. It is not a second independent circuit sample. All three initial observations were actually performed; their shared 37.988 s observation cost is separate from repair-controller times. A prior working handoff's unverified suggestion that observed R1/Tf commands differed is superseded by the exact identity receipt; no paper result relies on that suggestion.

\paragraph{Immutable graph data removes repeated XML scanning from candidate audits.}
The engineering variant builds a read-only cache from the actual RR XML, binding its digest, all payload digests, node metadata/capacities, directed edge--switch triples and native-rounded switch delays. It rejects nonzero RC rather than substituting a delay model. Every cache load rechecks the actual XML digest and all payload digests. Candidate scope, retained edges, terminal sets, full-tree reconstruction and current occupancy are still checked; native FULL/STA and the separate query audits remain unchanged. The original XML auditor remains frozen as the reference. Only candidate physical-audit data access changes in controller v4.

The first cache build covers 1028568 nodes and 4994132 directed edges in 33.122 s. Cached audits of the previously accepted Tf joint candidate and rejected R1 successor reproduce the complete route files byte for byte, as well as scope, query, attempt and changed-delay records. Their complete new audit child processes take 4.873 and 4.872 s, including digest verification and output writes; cache load/digest segments take 0.644 and 0.640 s. Their historical XML audit segments were 28.286 and 27.729 s, which motivate the optimization but are not a contemporaneous formal speedup comparison. A third accepted SPMV joint candidate also reproduces identical route bytes and audit fields; only its 3.270 s internal audit segment was measured, so no complete child time is inferred for it. Positive valid-edge and exact-capacity controls pass; six negative controls reject a nonexistent in-range edge, invalid switch, invalid node, over-capacity occupancy, different XML digest and mismatched payload digest. No routing or STA experiment is repeated for these engineering checks.

The next separately registered SPMV experiment starts from unmodified A0 with both exact policies. A 0.05 ns selection window leaves the same six worst connections, whereas 0.2 ns selects seventeen. That development window is used to test strict-reject/endpoint-accept plateau repairs and any later CPD benefit, with fresh target selection after each acceptance. It also exercises the candidate audit cache in the live automatic flow. Its unfinished outcomes are excluded from the present claims; neither worst-only exhaustion nor the cache validation ends the research.
